\chapter{Обучение на дофамине у червя и мухи}
\chaptermark{Червь и муха}
\label{sec:wormfly}

\section{Одно вещество, три схемы}

В главе~\ref{sec:dofa} \nt{DOPA}-нейроны рассылали по мозгу одно число --- ошибку $\delta$, а каждый синапс по своему следу решал, относится ли она к нему. Это решение для большой машины: у человека $8.6\cdot10^{10}$ нейронов~\cite{Azevedo2009}, и \nt{DOPA}-нейронов среди них ничтожная доля (таблица~\ref{tab:types}). Посмотрим на две машины поменьше, схема которых известна почти целиком. У червя $302$ нейрона, и их связи восстановлены по электронным микрофотографиям~\cite{White1986,Cook2019}. У мухи в мозге $139\,255$ нейронов и $5\cdot10^7$ синапсов между ними, и это тоже прочитано целиком~\cite{Dorkenwald2024}.

Дофамин есть у обоих, и оба учатся. Вопрос главы инженерный: что несёт провод \nt{DOPA}, кому он адресован и какое правило записи включает. Ответы получатся разные. У червя это флаг состояния, у мухи --- жгут адресных линий записи, каждая к своему банку памяти, у млекопитающих --- широковещательная ошибка главы~\ref{sec:dofa}.

\section{Червь: флаг «под нами еда»}

\emph{Аналоговая машина.} У изученных нейронов червя нет классических натриевых импульсов. Записанный сенсорный нейрон почти эквипотенциален и работает плавным потенциалом; чувствительность его падает с деполяризацией, и это расширяет диапазон входа~\cite{Goodman1998} --- регулировка усиления главы~\ref{sec:sensors}, сделанная калиевым током. Язык главы~\ref{sec:code} здесь другой: вместо счёта импульсов --- уровень напряжения.

\emph{Провод \nt{DOPA} --- выход датчика.} \nt{DOPA}-нейронов у червя восемь, и все они, по-видимому, механочувствительные~\cite{Sulston1975}. Они чувствуют механическое свойство бактерий, которыми червь питается: сытый червь на бактериях ползёт медленнее, и это замедление идёт через \nt{DOPA}~\cite{Sawin2000}. Значит, провод \nt{DOPA} у червя ничего не вычисляет: он сообщает, есть ли под телом еда. Это флаг, а не ошибка.

\emph{Кому адресован флаг.} Дофамин действует на мотонейроны, к которым от \nt{DOPA}-нейронов синапсов нет, --- вне синапсов, по объёму~\cite{Chase2004}. На одних и тех же мотонейронах стоят рецепторы двух семейств с противоположным действием, и поведение задаёт их разность. Знак сигнала выбирает получатель (утверждение~\ref{prop:receptor}), а у одного получателя знак складывается из двух рецепторов.

\emph{Что меняет флаг.} Во-первых, режим движения: после ухода с еды червь часто поворачивает, и частота поворотов спадает со временем после последней встречи с едой; так поиск остаётся около места, где еда была. Без \nt{DOPA}-нейронов этот поиск исчезает~\cite{Hills2004}. Для инженера это регистр, который заряжается от датчика и разряжается, как RC-цепь: память о недавней еде без единого изменения весов.

Во-вторых, скорость обучения. Червь, которого раз за разом встряхивают постукиванием, отвечает всё слабее --- привыкает. Без еды он привыкает быстрее, чем на еде. Сведения о еде поставляют механочувствительные \nt{DOPA}-нейроны, а дофамин через рецептор D1-типа меняет ответы датчиков прикосновения~\cite{Kindt2007,Sanyal2004}. Датчики прикосновения сами возбуждают \nt{DOPA}-нейроны, так что в схеме есть и положительная обратная связь. Запишем привыкание как спад ответа $A_k$ на $k$-е постукивание:
\begin{equation}
A_{k+1} \;=\; A_k \;-\; \eta(D)\,\bigl(A_k - A_\infty\bigr),
\label{eq:habituation}
\end{equation}
где $A_\infty$ --- уровень, к которому ответ стремится, а $D$ --- флаг еды. Флаг меняет скорость $\eta$, но не направление изменения и не то, какие синапсы меняются. Это та же роль, что у \nt{ACET} главы~\ref{sec:acet}: один уровень, общий для всех получателей, переключает режим записи.

В-третьих, ассоциации. Червь, долго пробывший в соли без еды или вместе с отталкивающим веществом, начинает соль избегать; в этом участвуют дофамин и серотонин~\cite{Hukema2008}. Для ассоциативного обучения червю нужен канал в пресинаптических окончаниях \nt{DOPA}-нейронов, который усиливает дофаминовую передачу~\cite{Voglis2008}. Несёт ли провод червя ошибку предсказания, как у млекопитающих, не показано.

\section{Муха: хэш и банки памяти}

Обучение мухи изучено лучше всего на запахах. Муха связывает запах с сахаром или с ударом тока и потом летит к запаху или от него. Схема, в которой это происходит, --- грибовидное тело --- устроена как запоминающее устройство инженерного учебника (рис.~\ref{fig:mushroom}).

\begin{figure}[t]
\centering
\begin{tikzpicture}[>=Stealth,
  kc/.style={circle,draw,thick,fill=blue!6,minimum size=0.3cm,inner sep=0pt},
  mbon/.style={circle,draw,very thick,fill=red!10,minimum size=0.75cm,inner sep=0pt,font=\scriptsize},
  dan/.style={draw,thick,rounded corners=2pt,fill=orange!15,font=\scriptsize,inner sep=2pt,align=center},
  lab/.style={font=\scriptsize,align=center}]
% input channels
\foreach \i in {0,...,5}{\draw[->,thick,blue!60!black] (-0.6,{0.35*\i}) -- (0.25,{0.35*\i});}
\node[lab,left] at (-0.65,0.9) {$\sim50$\\каналов\\запаха};
% expander
\foreach \i in {0,...,9}{\node[kc] (k\i) at ({0.75+0.3*mod(\i,2)},{-0.35+0.25*\i}) {};}
\foreach \i in {0,...,9}{\pgfmathsetmacro{\a}{0.35*mod(3*\i,6)}\pgfmathsetmacro{\b}{0.35*mod(3*\i+2,6)}
  \draw[gray!60!black] (0.25,\a) -- (k\i); \draw[gray!60!black] (0.25,\b) -- (k\i);}
\node[lab] at (0.9,-0.9) {расширитель\\$\sim2000$ \nt{GLUT},\\по $\sim7$ входов};
% feedback inhibition
\node[circle,draw,thick,fill=red!8,minimum size=0.55cm,font=\tiny] (apl) at (0.9,2.95) {\nt{GABA}};
\draw[->,thick,red!70!black] (1.3,2.2) -- (apl);
\draw[-{Bar[width=6pt]},thick,red!70!black] (apl.west) -- ++(-0.35,0) |- (0.62,2.0);
\node[lab,right] at (1.25,3.05) {держит активными $5$--$10\%$};
% axons through compartments
\foreach \y in {0.2,0.6,1.0,1.4}{\draw[thick,blue!60!black] (1.4,\y) -- (9.6,\y);}
\foreach \x/\n/\s in {3.3/1/{к запаху},5.9/2/{от запаха},8.5/15/{к запаху}}{
  \draw[thick,densely dashed,gray!60!black,rounded corners=3pt] (\x-1.0,-0.05) rectangle (\x+1.0,1.65);
  \node[mbon] (m\n) at (\x,-0.85) {выход};
  \foreach \y in {0.2,0.6,1.0,1.4}{\draw[gray!60!black] (\x-0.25,\y) -- (m\n);}
  \node[lab,below] at (\x,-1.25) {голос: \s};
  \node[dan] (d\n) at (\x,2.55) {учитель $c=\n$};
  \draw[->,very thick,orange!85!black] (d\n) -- (\x,1.7);}
\node[font=\small] at (7.2,0.8) {$\cdots$};
\node[lab] at (5.9,3.3) {$20$ типов \nt{DOPA}-нейронов, каждый пишет в свой отсек (из $15$)};
\end{tikzpicture}
\caption{Грибовидное тело мухи как запоминающее устройство, схема. Около пятидесяти каналов запаха случайно смешиваются в $\sim2000$ клетках расширителя; один \nt{GABA}-нейрон обратной связи оставляет активной малую долю клеток --- это адрес запаха. Аксоны расширителя проходят через $15$ отсеков; в каждом они подключены к своим выходам, и в каждый отсек пишет свой учитель --- \nt{DOPA}-нейроны одного типа. Выходы голосуют «к запаху» или «от запаха».}
\label{fig:mushroom}
\end{figure}

\emph{Адрес --- случайный хэш.} Запах приходит по пятидесяти с небольшим входным каналам~\cite{Caron2013}. Дальше стоит расширитель: около $2000$ \nt{GLUT}-клеток~\cite{Aso2014}, и каждая суммирует от $2$ до $11$ каналов, в среднем семь, выбранных, насколько удалось установить, случайно~\cite{Caron2013}. На данный запах отвечают лишь $5$--$10\%$ клеток. Такую разреженность держит один \nt{GABA}-нейрон: клетки расширителя возбуждают его, а он тормозит их всех --- чем больше горит, тем сильнее торможение~\cite{Lin2014}. Это выбор нескольких сильнейших через торможение, как в главе~\ref{sec:gaba}; назовём его отбором $k$ победителей. Если разорвать эту петлю, код становится плотнее, коды разных запахов --- похожее, и муха перестаёт различать близкие запахи, хотя далёкие различает~\cite{Lin2014}. Доля отвечающих клеток мало зависит и от концентрации запаха; в опыте в этом участвует привыкание к недавним запахам~\cite{Honegger2011}, а у отбора $k$ победителей это свойство встроено: порядок сумм не меняется, когда все входы растут вдвое.

Инженер узнает здесь алгоритм поиска похожего: случайная проекция в пространство в сорок раз большей размерности и отбор нескольких самых больших координат. Похожие запахи получают перекрывающиеся адреса, непохожие --- почти непересекающиеся~\cite{Dasgupta2017}. Тот же приём мы видели в расширителе цепи обучения движениям (глава~\ref{sec:motorlearning}), и в обоих случаях число входов у клетки расширителя близко к тому, при котором размерность кода наибольшая~\cite{LitwinKumar2017}.

\emph{Банки с отдельными линиями записи.} Аксоны расширителя проходят через $15$ отсеков. В каждом отсеке они подключены к выходным нейронам своего типа --- всего $34$ выходных нейрона $21$ типа, --- а \nt{DOPA}-нейроны $20$ типов посылают аксоны каждый в один, самое большее в два отсека~\cite{Aso2014,AsoMBON2014}. В терминах памяти: адрес --- разреженный код запаха, банк --- отсек, данные --- веса синапсов расширителя на выход отсека, линия разрешения записи --- \nt{DOPA}-нейроны этого отсека. Выходы голосуют: одни при возбуждении привлекают муху, другие отталкивают, и поведение задаёт весь их набор~\cite{AsoMBON2014}. Утверждение~\ref{prop:bundle} о жгуте здесь доведено до предела: у каждой линии свой адресат.

\emph{Учителя.} Двенадцать \nt{DOPA}-нейронов одной группы сообщают о наказании: если возбуждать их светом, пока муха нюхает запах, она потом избегает запаха так же, как после удара током~\cite{ClaridgeChang2009}. Другая группа сообщает о награде: эти нейроны возбуждаются, когда муха ест сахар, сильнее у голодной, и без них сахар перестаёт учить, хотя рефлекторный ответ на сахар остаётся~\cite{Liu2012}.

\emph{Правило записи.} Если запах сочетать с возбуждением \nt{DOPA}-нейронов отсека, синапсы расширителя на выход этого отсека ослабевают, и только для этого запаха. Изменение требует правильного порядка --- сначала запах, потом дофамин --- с точностью до долей секунды и остаётся в пределах отсека. Оно идёт и тогда, когда импульсы выходного нейрона полностью подавлены~\cite{Hige2015}: импульсы получателя не нужны, хватает двух множителей --- активности клетки расширителя и учителя. Пример того, как это меняет поведение: выход одного из отсеков при возбуждении привлекает муху, а при обучении с наказанием его ответ на этот запах падает~\cite{Hige2015}. Учитель наказания приглушает голос «к запаху»; что так устроены все отсеки, --- гипотеза.

Совпадение двух сигналов опознаёт одна молекула. Фермент, производящий внутриклеточный сигнал цАМФ, включается и кальцием, который входит в активную клетку, и рецептором дофамина; вместе они дают больше суммы, но только если деполяризация предшествует дофамину, и с теми же интервалами, что нужны для обучения~\cite{Tomchik2009,Levin1992}. Это вентиль И с памятью о порядке, сделанный из одного белка.

Порядок задаёт и знак. Дофамин после запаха ослабляет синапсы через рецептор одного типа, дофамин до запаха усиливает их через рецептор другого, и один и тот же \nt{DOPA}-нейрон может записать и влечение, и отвращение~\cite{Handler2019,AsoRubin2016}. В поведении это видно давно: удар током \emph{после} запаха учит его избегать, а удар \emph{до} запаха делает запах привлекательным --- он предвещает конец неприятности~\cite{Tanimoto2004}. Сведём всё в одно правило для веса $w_{cj}$ от клетки расширителя $j$ к выходу отсека $c$:
\begin{empheq}[box=\keybox]{equation}
\frac{\dd w_{cj}}{\dd t} \;=\; -\,\eta_c\,e_j\,d_c \;+\; \eta'_c\,\bar d_c\,x_j \;+\; \frac{w_0-w_{cj}}{\tau_c},
\qquad
\tau_e\frac{\dd e_j}{\dd t}=-e_j+x_j, \quad
\tau_d\frac{\dd \bar d_c}{\dd t}=-\bar d_c+d_c .
\label{eq:mbrule}
\end{empheq}
Здесь $x_j$ --- активность клетки $j$, $d_c$ --- активность учителя отсека $c$, $e_j$ и $\bar d_c$ --- их следы. Первый член: дофамин пришёл, пока жив след запаха, --- ослабление. Второй: запах пришёл, пока жив след дофамина, --- усиление. Третий: вес сам возвращается к исходному с постоянной времени отсека $\tau_c$. Сравните с~\eqref{eq:threefactor}: там след $e=xy$ требовал активности получателя, а сигнал $\delta$ был один на всех; здесь след принадлежит отправителю, а у каждого банка своя линия $d_c$ и свои $\eta_c$, $\tau_c$.

\emph{Банки разной скорости.} Отсеки пишут и забывают по-разному. Один учитель наказания записывает прочную память за одно сочетание, но через сутки она почти стёрта, и четырёхдневной памяти он не создаёт. Другой после одного сочетания почти ничего не оставляет, зато десять сочетаний с перерывами дают память на четыре дня. В целом быстрая запись и долгое хранение, по-видимому, связаны обратно~\cite{AsoRubin2016}; для ослабления синапсов в одном отсеке нужно гораздо более долгое сочетание, чем в другом~\cite{Hige2015}. В~\eqref{eq:mbrule} это разные $\eta_c$ и $\tau_c$: кэш и диск в одной схеме.

\emph{Откуда учитель знает ожидание.} Выходные нейроны отсеков посылают связи к \nt{DOPA}-нейронам~\cite{Li2020}. Через такие петли муха переоценивает выученное: если запах, обещавший сахар, приходит без сахара, выходы возбуждают учителей противоположного знака, и новая память гасит старую; если обещание сбылось, петля подтверждает прежнюю память~\cite{Felsenberg2017}. Так учитель получает на вход и подкрепление, и предсказание, сделанное выходами, --- то, что нужно для ошибки~\eqref{eq:ctd}. Что сигнал учителей мухи в целом и есть ошибка предсказания, --- гипотеза.

\begin{keyblock}
\begin{proposition}[Адрес учителя]
\label{prop:wormfly}
Провод \nt{DOPA} может нести флаг режима без адреса --- он меняет скорость и усиление у всех получателей сразу (червь); набор линий записи, каждая к своему банку, где адрес данных --- разреженный код входа, а знак записи задаёт порядок событий (муха, правило~\eqref{eq:mbrule}); или одно широковещательное число ошибки, которое каждый синапс относит к себе по своему следу (глава~\ref{sec:dofa}). Чем подробнее адрес у линии учителя, тем меньше синапсу нужно знать самому: у мухи хватает двух множителей, в~\eqref{eq:threefactor} их три. Устройство и правила записи грибовидного тела измерены; что это общий закон для машин разного размера, --- гипотеза.
\end{proposition}
\end{keyblock}

\section{Модель: хэш и один банк}

Проверим, что даёт разреженный адрес, на модели \texttt{mushroom\_body.py}. Запах --- $50$ случайных неотрицательных чисел, по одному на канал. Каждая из $2000$ клеток расширителя суммирует $7$ случайных каналов, торможение оставляет активной долю $f$ клеток с самыми большими суммами. Один выход с весами $w_j=1$ голосует «к запаху». Одно сочетание запаха $A$ с учителем наказания обнуляет веса всех клеток, активных для $A$. Отвращение к запаху $X$ --- доля его активных клеток, чьи веса обнулены.

Результат (рис.~\ref{fig:mbmodel}, двадцать моделей-мух). Выученный запах $A$ муха избегает полностью. Запах $A'$, в котором на $10$, $20$ и $30\%$ подмешан другой, при $f=5\%$ она избегает на $0.91$, $0.81$ и $0.72$, а совсем другой запах $B$ --- на $0.05$. Если ослабить торможение до $f=20\%$, различение хуже: $0.94$, $0.87$, $0.81$, и непохожий запах $B$ избегается уже на $0.20$. При $f=50\%$ муха избегает половины всех запахов после одного урока. Чем реже код, тем точнее адрес, --- то же, что показал разрыв петли торможения у настоящей мухи~\cite{Lin2014}. Шума в модели нет, поэтому нижнего предела разреженности она не показывает: при слишком редком коде ответ держится на немногих клетках и становится ненадёжным (глава~\ref{sec:code}).

\begin{figure}[t]
\centering
\begin{tikzpicture}[x=1cm,y=3.2cm]
\draw[->,gray!70!black] (0,0) -- (11.2,0);
\draw[->,gray!70!black] (0,0) -- (0,1.12) node[above,font=\small,black] {отвращение};
\foreach \v in {0.25,0.5,0.75,1.0}{\draw[gray!60!black] (-0.06,\v) -- (0.06,\v); \node[left,font=\scriptsize] at (0,\v) {\v};}
\foreach \x/\a/\b/\lab in {1.1/1.00/1.00/{$A$},3.2/0.91/0.94/{$A'$, $10\%$},5.3/0.81/0.87/{$A'$, $20\%$},7.4/0.72/0.81/{$A'$, $30\%$},9.5/0.05/0.20/{$B$}}{
  \fill[blue!45] (\x-0.75,0) rectangle (\x-0.05,\a);
  \fill[red!55!black] (\x+0.05,0) rectangle (\x+0.75,\b);
  \node[below,font=\scriptsize] at (\x,-0.01) {\lab};
  \node[above,font=\scriptsize] at (\x-0.4,\a) {\a};
  \node[above,font=\scriptsize] at (\x+0.4,\b) {\b};}
\fill[blue!45] (7.9,1.07) rectangle (8.2,1.0); \node[right,font=\scriptsize] at (8.25,1.035) {$f=5\%$};
\fill[red!55!black] (9.5,1.07) rectangle (9.8,1.0); \node[right,font=\scriptsize] at (9.85,1.035) {$f=20\%$};
\end{tikzpicture}
\caption{Модель \texttt{mushroom\_body.py}: насколько муха избегает запаха после одного урока с запахом $A$. $A'$ --- запах $A$ с подмешанным на $10$--$30\%$ другим, $B$ --- непохожий запах. Среднее по двадцати случайным схемам. Синие столбцы --- разреженный код ($5\%$ активных клеток), тёмно-красные --- ослабленное торможение ($20\%$).}
\label{fig:mbmodel}
\end{figure}

\section{Итог}

\begin{table}[t]
\centering
\footnotesize
\setlength{\tabcolsep}{3pt}
\renewcommand{\arraystretch}{1.15}
\begin{tabular}{>{\raggedright}p{2.0cm}>{\raggedright}p{3.4cm}>{\raggedright}p{3.9cm}>{\raggedright\arraybackslash}p{3.6cm}}
\toprule
& Червь & Муха & Млекопитающие (гл.~\ref{sec:dofa}) \\
\midrule
нейронов & $302$ & $1.4\cdot10^5$ & $8.6\cdot10^{10}$ у человека \\
\nt{DOPA}-нейроны & $8$, механочувствительные & $20$ типов у грибовидного тела & доля по табл.~\ref{tab:types} \\
что несёт провод & флаг «под нами еда» & награду или наказание для своего банка & ошибку предсказания $\delta$ \\
адрес & нет: объём, знак выбирает рецептор & отсек-банк; данные адресует разреженный код & нет: адрес --- след в синапсе \\
правило & скорость привыкания~\eqref{eq:habituation}, режим движения & два множителя, знак по порядку~\eqref{eq:mbrule}, свои $\eta_c$, $\tau_c$ & три множителя~\eqref{eq:threefactor} \\
что известно & роль контекста измерена; ошибка предсказания не показана & схема и правило записи измерены; ошибка предсказания --- гипотеза & сходство с $\delta$ измерено \\
\bottomrule
\end{tabular}
\caption{Три схемы обучения на одном веществе.}
\label{tab:wormfly}
\end{table}

Одна молекула служит трём разным схемам (таблица~\ref{tab:wormfly}). У червя \nt{DOPA} --- датчик, выведенный на общий провод: флаг еды меняет режим движения и скорость привыкания всех получателей сразу. У мухи он разложен на двадцать адресных линий записи, и память устроена как банк с хэш-адресацией: разреженный код запаха выбирает строки, учитель отсека разрешает запись, порядок событий задаёт её знак, а у каждого банка своя скорость записи и забывания. У млекопитающих линия снова одна, но несёт уже вычисленную ошибку, а адрес держит след в каждом синапсе. Для инженера главный урок --- в утверждении~\ref{prop:wormfly}: сколько знаний о том, куда писать, положено в провод учителя, столько не нужно класть в синапс.

\section{Задачи}

\begin{exercise}[\lvlA]
\label{ex:23:1}
\emph{Сколько запахов в банке.} В одном отсеке каждое наказание обнуляет веса всех клеток, активных для запаха; доля активных клеток~$f$. После скольких наказанных запахов ответ отсека на новый, никогда не встречавшийся запах упадёт вдвое при $f=5\%$, $10\%$, $50\%$? Зачем тогда забывание $\tau_c$ в~\eqref{eq:mbrule}?
\end{exercise}

\begin{exercise}[\lvlA]
\label{ex:23:2}
\emph{Семь проводов из пятидесяти одного.} Клетка расширителя берёт $K$ случайных каналов из $51$. При каком наименьшем $K$ среди $2000$ клеток в среднем меньше одной пары с одинаковым набором? Какую долю входов в среднем делят две клетки при $K=7$ и $K=25$?
\end{exercise}

\begin{exercise}[\lvlA]
\label{ex:23:3}
\emph{Привыкание на еде.} Червя постукивают раз в $10$~с; по~\eqref{eq:habituation} $\eta=0.1$ без еды и $\eta=0.05$ на еде. Через сколько постукиваний и секунд превышение ответа над $A_\infty$ падает вдвое в каждом случае? Меняет ли флаг конечный уровень ответа?
\end{exercise}

\begin{exercise}[\lvlB]
\label{ex:23:4}
\emph{Знак в середине запаха.} В~\eqref{eq:mbrule} пусть $\tau_e=\tau_d=\tau$, $\eta'_c=\eta_c$, забыванием пренебрежём; запах $x_j=1$ идёт на $[0,T]$, а учитель даёт короткий импульс $d_c=\tau\,\delta(t-\Delta)$. Найдите изменение веса как функцию $\Delta$ и докажите, что знак меняется ровно при $\Delta=T/2$, каковы бы ни были $T$ и $\tau$.
\end{exercise}

\begin{exercise}[\lvlB]
\label{ex:23:5}
\emph{Быстрый банк и медленный.} Банк 1: $\eta=0.9$, $\tau=0.5$~сут; банк 2: $\eta=0.05$, $\tau=10$~сут; сочетание даёт $m\leftarrow m+\eta(1-m)$, потом $m$ спадает с $\tau$. Найдите сумму двух банков через минуту, сутки и четыре дня после одного сочетания и после десяти подряд. Какой банк держит память на четвёртый день?
\end{exercise}

\begin{exercise}[\lvlB]
\label{ex:23:6}
\emph{Громкий запах.} В \texttt{mushroom\_body.py} замените торможение неподвижным порогом, при котором запах $A$ зажигает $5\%$ клеток, и предъявите $A$ вдвое сильнее. Какая доля клеток загорится и насколько муха будет избегать выученного запаха? А с отбором $k$ победителей?
\end{exercise}

\begin{exercise}[\lvlC]
\label{ex:23:7}
\emph{Критик из выходов.} Постройте из двух отсеков (награда и наказание) и связей выходов на учителей схему, в которой запах, обещавший сахар и пришедший без сахара, стирает память, а сбывшееся обещание её не трогает~\cite{Felsenberg2017}. Какой знак нужен у связей и какую величину~\eqref{eq:ctd} несут тогда учителя?
\end{exercise}
